\documentclass[12pt,a4paper]{article}

\usepackage[utf8]{inputenc}
\usepackage[T1]{fontenc}
\usepackage[polish]{babel}
\usepackage{lmodern}
\usepackage[a4paper,margin=2.5cm]{geometry}
\usepackage{amsmath}
\usepackage{graphicx}
\usepackage{enumitem}
\usepackage{booktabs}
\usepackage{array}
\usepackage{tabularx}
\usepackage[hidelinks]{hyperref}

\setlist[itemize]{topsep=0.35em,itemsep=0.2em,leftmargin=1.6em}
\setlist[enumerate]{topsep=0.35em,itemsep=0.2em,leftmargin=1.6em}

\newcommand{\kod}[1]{\texttt{#1}}
\setlength{\emergencystretch}{3em}
\newcolumntype{L}{>{\raggedright\arraybackslash}X}

\begin{document}

% Nagłówek jak we wzorze materialy/zal_proj_zssk.pdf (bez strony tytułowej).
\thispagestyle{empty}
\noindent
\begin{minipage}[t]{0.5\textwidth}
	Karol Gąsior\\
	Indeks 280913\\[0.6em]
	Mateusz Kunicki\\
	Indeks 272468\\[0.6em]
	Grupa G2 (pon. 13:40)
\end{minipage}%
\begin{minipage}[t]{0.5\textwidth}
	\raggedleft Wrocław, 09-11-2026 r.
\end{minipage}

\vspace{1.5em}
\begin{center}
	{\Large\bfseries Opis założeń projektu}
\end{center}
\vspace{0.5em}

%======================================================================
\section{Temat i cel projektu}

Temat: Mechanizmy zarządzania wątkami programowymi w zadaniu wykrywania krawędzi na obrazach cyfrowych z wykorzystaniem operatorów gradientowych (Sobela i pokrewnych).

Cel projektu: projekt i implementacja programu wykrywającego krawędzie na obrazie cyfrowym w pięciu wersjach oraz zmierzenie, ile czasu przetwarzania oszczędza każdy rodzaj przyspieszenia. Wersja W1 jest prostą implementacją sekwencyjną i stanowi punkt odniesienia. Każda z wersji W2--W5 bierze kod W1 i wprowadza do niego \emph{jeden} rodzaj przyspieszenia, odpowiadający innemu poziomowi przetwarzania równoległego \cite{czech}: zmianę algorytmu i sposobu korzystania z pamięci (W2), równoległość na poziomie instrukcji procesora -- SIMD (W3), równoległość na poziomie wątków na wielu rdzeniach procesora (W4) oraz przetwarzanie na setkach rdzeni karty graficznej (W5). Wersje nie są łączone ze sobą: każda jest porównywana z W1, dzięki czemu wynik pokazuje wkład jednego mechanizmu, a nie sumę kilku.

Zakres projektu obejmuje:
\begin{itemize}
	\item maski Sobela $3\times3$, $5\times5$ i $7\times7$, Prewitta, Scharra i Robertsa;
	\item obrazy od 0,04 do 59~Mpix (zdjęcia i obrazy syntetyczne);
	\item w wersji wielowątkowej W4 -- mechanizmy zarządzania wątkami: liczbę wątków od 1 do 24 (oraz 32 i 48), cztery strategie przydziału pracy, rozmiar przydzielanego fragmentu, sposób synchronizacji wątków przy wyznaczaniu wartości maksymalnej, pulę wątków w porównaniu z tworzeniem wątków przy każdym etapie i przypisanie wątków do rdzeni procesora;
	\item w wersji GPU W5 -- liczbę wątków uruchamianych na karcie graficznej, rozmiar grupy roboczej i koszt przesyłania danych między pamięcią komputera a kartą;
	\item dwa pomiary każdej wersji: samego algorytmu oraz całego programu z wczytaniem i zapisem obrazu;
	\item ocenę, czy większa maska daje lepszy wynik wykrywania krawędzi, mierzoną na obrazach o znanym położeniu krawędzi z dodanym szumem;
	\item propozycje dalszych przyspieszeń przygotowane z pomocą modelu sztucznej inteligencji.
\end{itemize}

%======================================================================
\section{Sformułowanie problemu}

Problem będzie rozwiązywany przy następujących założeniach:
\begin{itemize}
	\item obraz wejściowy jest wczytywany z pliku (JPG, PNG, GIF) i konwertowany do skali szarości o 8~bitach na piksel albo generowany syntetycznie w pamięci;
	\item cały obraz mieści się w pamięci operacyjnej;
	\item piksele leżące poza brzegiem obrazu zastępuje się najbliższym pikselem brzegowym (powielenie krawędzi);
	\item wynikiem jest obraz 8-bitowy o tych samych wymiarach co wejście;
	\item najmniejszą jednostką podziału pracy pomiędzy wątki procesora jest wiersz obrazu, a pomiędzy wątki karty graficznej -- piksel;
	\item wszystkie wersje muszą dawać wynik identyczny co do bitu z wynikiem wersji W1.
\end{itemize}

Problem obliczeniowy \cite{gonzalez,tadeusiewicz}: dla obrazu $I$ o wymiarach $W \times H$ i pary masek $K_x$, $K_y$ o rozmiarze $k \times k$ wyznaczyć dla każdego piksela moduł gradientu
\begin{equation}
	G(x, y) = \sqrt{G_x(x, y)^{2} + G_y(x, y)^{2}}, \qquad G_x = K_x * I, \qquad G_y = K_y * I,
	\label{eq:gradient}
\end{equation}
gdzie $*$ oznacza dwuwymiarowy splot dyskretny
\begin{equation}
	(K * I)(x, y) = \sum_{i=0}^{k-1} \sum_{j=0}^{k-1} K(i, j)\; I\!\left(x + j - \lfloor k/2 \rfloor,\; y + i - \lfloor k/2 \rfloor\right),
\end{equation}
a następnie przeskalować wynik do zakresu 0--255: $O(x, y) = \operatorname{round}\big(255 \cdot G(x, y) / \max G\big)$. Skalowanie względem największej wartości w obrazie wymaga znajomości $\max G$, czyli zakończenia obliczeń dla całego obrazu przed zapisaniem pierwszego piksela wyniku. Jako warianty pomocnicze zostaną zbadane skalowanie stałe (względem teoretycznego maksimum maski) i obcięcie wartości większych od 255, które tego wymagania nie mają.

Problem decyzyjny: czy dla danego obrazu, maski i sprzętu istnieje wersja programu i jej konfiguracja (liczba wątków, strategia przydziału pracy, rozmiar fragmentu, rozmiar grupy roboczej GPU), dla której czas przetwarzania nie przekracza zadanej wartości $T_{\max}$?

Problem optymalizacyjny \cite{williams,openmp}: dla każdej wersji dobrać parametry tak, aby zapewnić optymalne wartości kryteriów:
\begin{itemize}
	\item czas zegarowy przetwarzania obrazu -- minimum;
	\item przyspieszenie względem wersji W1, $S = T_{W1} / T_{Wi}$ -- maksimum;
	\item dla wersji wielowątkowej: przyspieszenie $S(N) = T(1)/T(N)$ i efektywność $E(N) = S(N)/N$ -- maksimum; nierównowaga obciążenia wątków (różnica czasów zakończenia pracy) -- minimum;
	\item zgodność wyniku z wersją W1 -- wynik identyczny co do bitu.
\end{itemize}
Niezależnym kryterium jakości jest trafność wykrywania krawędzi (rozdział~\ref{sec:testy}), która zależy od maski, ale nie od wersji programu.

%======================================================================
\section{Analiza złożoności obliczeniowej problemu}

Wyznaczenie modułu gradientu jest problemem z klasy P. Splot dwuwymiarowy wymaga dla każdego piksela $k^2$ mnożeń i dodawań dla $G_x$ i tyle samo dla $G_y$, czyli $2k^2$ operacji, a więc całe zadanie ma złożoność $O(W H k^{2})$. Maski Sobela, Prewitta i Scharra są rozdzielne -- dają się zapisać jako iloczyn wektora wygładzającego $s$ i różnicującego $d$, np. dla Sobela $3\times3$:
\begin{equation}
	K_x = s^{T} d = \begin{bmatrix} 1 \\ 2 \\ 1 \end{bmatrix} \begin{bmatrix} -1 & 0 & 1 \end{bmatrix},
	\qquad K_y = d^{T} s.
\end{equation}
Splot z maską rozdzielną można wtedy policzyć jako dwa sploty jednowymiarowe: poziomy (dla każdego wiersza obrazu wektorami $s$ i $d$) i pionowy (na wynikach przebiegu poziomego). Wymaga to $4k$ zamiast $2k^2$ operacji na piksel, co daje złożoność $O(W H k)$: dla $k = 3$ jest to 12 zamiast 18 operacji, a dla $k = 7$ -- 28 zamiast 98. Maska Robertsa nie jest rozdzielna. Skalowanie względem maksimum wymaga dodatkowego przebiegu $O(W H)$ i redukcji maksimum. Złożoność pamięciowa wynosi $O(W H)$: obraz wejściowy (1~B/piksel), bufor modułu gradientu (4~B/piksel) i wynik (1~B/piksel), łącznie około 6~B/piksel.

Wartość każdego piksela wyniku zależy wyłącznie od jego otoczenia w obrazie wejściowym, więc zadanie jest łatwo zrównoleglane i teoretycznie pozwala na przyspieszenie liniowe względem liczby jednostek obliczeniowych. Przyspieszenie ogranicza jednak część sekwencyjna programu, zgodnie z prawem Amdahla \cite{amdahl,czech}:
\begin{equation}
	S(N) \le \frac{1}{s + \dfrac{1 - s}{N}}, \qquad \lim_{N \to \infty} S(N) = \frac{1}{s},
	\label{eq:amdahl}
\end{equation}
gdzie $s$ oznacza udział części sekwencyjnej w czasie wykonania na jednej jednostce. W tym zadaniu część sekwencyjną tworzą: wczytanie i dekodowanie pliku obrazu, zapis wyniku, alokacja buforów, redukcja maksimum i bariera między etapami, a w wersji GPU -- przesłanie obrazu na kartę i wyniku z powrotem. Drugim ograniczeniem jest przepustowość pamięci: splot $3 \times 3$ wykonuje kilkanaście operacji na każdy bajt danych, więc po przyspieszeniu obliczeń może go ograniczać szybkość odczytu i zapisu pamięci, a nie moc obliczeniowa.

Liczbę operacji i sposób zrównoleglenia w poszczególnych wersjach zestawiono w tabeli~\ref{tab:zlozonosc}.

\begin{table}[ht]
	\centering
	\caption{Złożoność i równoległość wersji programu (obraz $W \times H$, maska $k \times k$, $N$ wątków).}
	\label{tab:zlozonosc}
	\small
	\begin{tabularx}{\textwidth}{@{}lLLl@{}}
		\toprule
		Wersja & Operacje na piksel & Równoległość & Czas (rząd) \\
		\midrule
		W1 & $2k^2$ mnożeń, $2k^2$ obcięć współrzędnych & brak & $O(WHk^2)$ \\
		W2 & $4k$ mnożeń (rozdzielne), bez obcięć we wnętrzu & brak & $O(WHk)$ \\
		W3 & $2k^2$ mnożeń wykonywanych po 8 naraz & 8 pikseli na instrukcję (AVX2) & $O(WHk^2/8)$ \\
		W4 & $2k^2$ mnożeń, jak W1 & $N$ wątków na rdzeniach CPU & $O(WHk^2/N)$ + część sekwencyjna \\
		W5 & $2k^2$ mnożeń, jak W1 & tysiące wątków GPU na 768 rdzeniach & $O(WHk^2/P)$ + przesył danych \\
		\bottomrule
	\end{tabularx}
\end{table}

%======================================================================
\section{Metoda i algorytmy rozwiązywania problemu}

\subsection{Wersja W1 -- sekwencyjna (punkt odniesienia)}

Wersja W1 realizuje wzór~\eqref{eq:gradient} wprost, w jednym wątku. Dla każdego piksela $(x, y)$ obrazu program:
\begin{enumerate}
	\item przechodzi po wszystkich $k \times k$ polach maski; dla każdego pola wyznacza współrzędne piksela źródłowego i obcina je do zakresu obrazu funkcją \kod{std::clamp} (także wtedy, gdy piksel leży daleko od brzegu i obcięcie nie jest potrzebne);
	\item odczytuje współczynniki $K_x$ i $K_y$ z tablicy dynamicznej (\kod{std::vector}), której rozmiar jest znany dopiero w czasie wykonania, i sumuje iloczyny w zmiennych całkowitych $G_x$ i $G_y$;
	\item oblicza moduł $\sqrt{G_x^2 + G_y^2}$ w arytmetyce podwójnej precyzji i zapisuje go jako liczbę \kod{float} do bufora modułu, aktualizując przy tym bieżące maksimum.
\end{enumerate}
Po przetworzeniu całego obrazu drugi przebieg mnoży każdą wartość bufora przez $255/\max G$, zaokrągla ją i zapisuje jako bajt wyniku. Bufor modułu i bufor wyniku są przydzielane i zerowane przy każdym wywołaniu. Wersja ta jest najprostszą poprawną implementacją i jej wynik jest wzorcem, z którym porównywane są wszystkie pozostałe wersje.

Moduł liczony jest w podwójnej precyzji celowo: iloczyny $G_x^2$ i $G_y^2$ są wtedy dokładne, a pierwiastek poprawnie zaokrąglony, więc wynik nie zależy od kolejności działań ani od tego, czy kompilator połączy mnożenie z dodawaniem w jedną instrukcję. Pozwala to wymagać od każdej wersji wyniku identycznego co do bitu.

\subsection{Wersja W2 -- optymalizacja algorytmu i dostępu do pamięci (1 rdzeń)}

Wersja W2 działa na jednym rdzeniu i nie używa instrukcji wektorowych w sposób jawny. Wprowadza cztery zmiany w samym algorytmie i organizacji danych:
\begin{enumerate}
	\item \textbf{Obsługa brzegu poza główną pętlą.} Obcinanie współrzędnych jest potrzebne tylko w pasach szerokości $\lfloor k/2 \rfloor$ przy brzegach obrazu. Wnętrze obrazu jest przetwarzane osobną pętlą bez \kod{std::clamp}, a wskaźniki do $k$ wierszy źródłowych są wyznaczane raz na wiersz wyniku zamiast przy każdym odczycie. Znikają rozgałęzienia i przeliczanie indeksów w najczęściej wykonywanym fragmencie kodu.
	\item \textbf{Maska jako stała czasu kompilacji.} Rozmiar maski i jej współczynniki są parametrami szablonu C++ (\kod{constexpr}). Kompilator rozwija pętle po masce, pomija mnożenia przez zero (dla Sobela $3\times3$ trzy z dziewięciu współczynników każdej maski są zerowe) i zamienia mnożenia przez małe stałe na tańsze operacje.
	\item \textbf{Splot rozdzielny.} Dla masek rozdzielnych każdy wiersz obrazu jest najpierw splatany poziomo z wektorami $s$ i $d$, a wyniki $k$ ostatnich wierszy są przechowywane w buforze pierścieniowym. Splot pionowy korzysta z tego bufora. Każdy wiersz źródłowy jest przetwarzany poziomo tylko raz, a liczba operacji spada z $2k^2$ do $4k$ na piksel. Bufor $k$ wierszy mieści się w pamięci podręcznej L1/L2 procesora.
	\item \textbf{Mniejsze typy i bufory wielokrotnego użytku.} Wyniki przebiegu poziomego mieszczą się w 16 bitach (dla Sobela $7\times7$ co najwyżej $255 \cdot 64 = 16\,320$), więc bufor pierścieniowy przechowuje liczby \kod{int16} zamiast \kod{int32}, co zmniejsza o połowę ilość przesyłanych danych. Bufory są wyrównane do 64~B (rozmiar linii pamięci podręcznej) i przydzielane raz, a nie przy każdym wywołaniu, co usuwa koszt przydziału i pierwszego zapisu stron pamięci.
\end{enumerate}
Ponieważ zysk ze splotu rozdzielnego rośnie z rozmiarem maski, oczekuje się, że W2 będzie zyskiwać względem W1 tym więcej, im większe jest $k$.

\subsection{Wersja W3 -- wektoryzacja SIMD (1 rdzeń)}

Wersja W3 zachowuje algorytm W1 bez zmian: splot dwuwymiarowy $k \times k$, współczynniki maski odczytywane z tablicy w czasie wykonania (bez pomijania zer i bez rozdzielności), ten sam sposób liczenia modułu, przydział buforów przy każdym wywołaniu. Jedyną zmianą jest wykonywanie obliczeń dla 8 pikseli jednocześnie instrukcjami wektorowymi AVX2 \cite{intel}, czyli równoległość na poziomie instrukcji procesora:
\begin{itemize}
	\item dla każdego pola maski 8 sąsiednich pikseli (8~bajtów) jest wczytywanych jedną instrukcją i rozszerzanych do ośmiu liczb 32-bitowych w 256-bitowym rejestrze (\kod{\_mm256\_\allowbreak{}cvtepu8\_epi32});
	\item współczynnik maski jest powielany do wszystkich 8 pozycji rejestru, a mnożenie i dodawanie (\kod{\_mm256\_\allowbreak{}mullo\_epi32}, \kod{\_mm256\_\allowbreak{}add\_epi32}) wykonują się dla 8 pikseli w jednej instrukcji;
	\item moduł jest liczony wektorowo: pierwiastkiem z 8 liczb pojedynczej precyzji (\kod{\_mm256\_\allowbreak{}sqrt\_ps}), gdy $G_x^2 + G_y^2 < 2^{24}$ i wynik jest dokładny, a w przeciwnym razie 2~razy po 4~liczby podwójnej precyzji (\kod{\_mm256\_\allowbreak{}sqrt\_pd});
	\item drugi przebieg (skalowanie do 0--255) również przetwarza 8 pikseli na instrukcję.
\end{itemize}
Instrukcji wektorowej nie da się zastosować do pikseli brzegowych, dla których każdy z 8 pikseli wymagałby innego obcięcia współrzędnych. Dlatego kolumny brzegowe są liczone kodem W1, a w ścieżce wektorowej obcinany jest tylko numer wiersza źródłowego (raz na wiersz maski). Teoretyczna granica przyspieszenia wynosi 8; w praktyce będzie mniejsza z powodu kosztu rozszerzania bajtów do 32 bitów, części skalarnej przy brzegach i przepustowości pamięci.

\subsection{Wersja W4 -- wątki na wielu rdzeniach procesora}

Wersja W4 wykonuje kod W1 bez zmian (to samo jądro obliczeniowe), ale dzieli obraz pomiędzy $N$ wątków. Zastosowana zostanie dekompozycja danych: obraz jest dzielony na fragmenty złożone z kolejnych wierszy, a każdy wątek zapisuje wyniki do rozłącznego obszaru pamięci, więc same obliczenia nie wymagają blokad. Przetwarzanie przebiega w dwóch etapach rozdzielonych barierą:
\begin{enumerate}
	\item splot i wyznaczenie lokalnego maksimum przez każdy wątek;
	\item redukcja lokalnych maksimów do maksimum globalnego, a następnie równoległe skalowanie wyniku.
\end{enumerate}

Wątki są tworzone raz i utrzymywane w puli (\emph{thread pool}): między etapami czekają uśpione na zmiennej warunkowej (\kod{std::condition\_variable}) i są budzone, gdy pojawia się praca. Przydziałem fragmentów zarządza scheduler oferujący cztery strategie wzorowane na klauzuli \kod{schedule} standardu OpenMP \cite{openmp}:
\begin{itemize}
	\item \emph{static} -- obraz dzielony na $N$ równych bloków, po jednym na wątek; brak narzutu, ale wątek opóźniony przez system wydłuża cały etap;
	\item \emph{cyclic} -- fragmenty po $c$ wierszy przydzielane kolejnym wątkom po kolei (fragment $i$ dla wątku $i \bmod N$);
	\item \emph{dynamic} -- wolny wątek pobiera następny fragment ze wspólnego licznika atomowego; dobre równoważenie obciążenia kosztem operacji atomowej na każdy fragment;
	\item \emph{guided} -- jak \emph{dynamic}, ale rozmiar fragmentu maleje wraz z ilością pozostałej pracy (duże fragmenty na początku, małe na końcu).
\end{itemize}

W wersji W4 zostaną zbadane następujące mechanizmy zarządzania wątkami:
\begin{itemize}
	\item liczba wątków $N$ od 1 do 24 oraz 32 i 48 (więcej wątków niż rdzeni logicznych), z uwzględnieniem topologii procesora: 12 rdzeni fizycznych po 2 wątki sprzętowe (SMT), podzielonych na 2 bloki (CCD) po 6 rdzeni z osobną pamięcią podręczną L3;
	\item strategia przydziału i rozmiar fragmentu $c$ od 1 do 256 wierszy;
	\item sposób wyznaczania maksimum globalnego: lokalne maksimum każdego wątku i redukcja po barierze; porównanie każdego piksela ze wspólną zmienną atomową (\kod{std::atomic}, zapis operacją porównaj-i-zamień); aktualizacja wspólnej zmiennej pod blokadą \kod{std::mutex} dla każdego piksela; zapis bieżącego maksimum do własnej komórki wątku przy każdym pikselu -- z komórkami w osobnych liniach pamięci podręcznej albo w jednej linii, co wywołuje fałszywe współdzielenie (\emph{false sharing}) \cite{drepper};
	\item pula wątków w porównaniu z tworzeniem nowych wątków (\kod{std::thread}) w każdym etapie;
	\item przypisanie wątków do rdzeni logicznych (\kod{pthread\_setaffinity\_np}): bez przypisania (decyduje system), najpierw jeden blok rdzeni, naprzemiennie oba bloki, tylko jeden wątek na rdzeń fizyczny;
	\item koszt bariery wynikającej ze skalowania względem maksimum, w porównaniu ze skalowaniem stałym, które wymaga tylko jednego etapu.
\end{itemize}
Oczekuje się, że przyspieszenie W4 będzie ograniczone przez część sekwencyjną kodu W1 (przydział i zerowanie buforów wykonywane jednym wątkiem przed podziałem pracy), barierę między etapami i przepustowość pamięci.

\subsection{Wersja W5 -- karta graficzna (OpenCL)}

Wersja W5 przenosi algorytm W1 bez zmian na kartę graficzną NVIDIA GeForce GTX 1050 Ti (768 rdzeni CUDA w 6 multiprocesorach) przy użyciu standardu OpenCL \cite{opencl}. Każdy wątek karty (\emph{work-item}) oblicza jeden piksel: czyta $k \times k$ pikseli z pamięci globalnej karty, mnoży je przez współczynniki maski z pamięci stałej i oblicza moduł. Przebieg programu:
\begin{enumerate}
	\item przesłanie obrazu z pamięci komputera do pamięci karty przez magistralę PCIe;
	\item kernel splotu; w każdej grupie roboczej wątki wyznaczają wspólne maksimum w pamięci lokalnej, a jeden wątek grupy aktualizuje maksimum globalne operacją atomową;
	\item odczyt maksimum przez procesor i obliczenie skali;
	\item kernel skalowania do 0--255;
	\item przesłanie wyniku z powrotem do pamięci komputera.
\end{enumerate}
Czas każdego etapu zostanie zmierzony osobno (zdarzenia profilujące OpenCL). Inicjalizacja (wybór urządzenia, kompilacja kernela) i przydział buforów na karcie odbędą się poza pomiarem.

Liczba rdzeni karty jest stała, ale programista wybiera liczbę wątków uruchamianych na karcie oraz rozmiar grupy roboczej (wątków współdzielących pamięć lokalną i wykonywanych na jednym multiprocesorze). Zostanie zbadany wpływ rozmiaru grupy roboczej ($8\times8$, $16\times16$, $32\times8$, $32\times32$) i liczby pikseli liczonych przez jeden wątek (1--64, czyli od kilku milionów do kilkuset tysięcy wątków dla obrazu 4K). Karta tej klasy liczy w podwójnej precyzji około 32 razy wolniej niż w pojedynczej, dlatego moduł będzie liczony w podwójnej precyzji tylko wtedy, gdy pojedyncza nie zapewnia wyniku dokładnego; dodatkowo zostanie zmierzony wariant zawsze w pojedynczej precyzji z oceną różnicy wyniku. Oczekuje się, że przesył danych przez PCIe będzie stanowił część sekwencyjną ograniczającą zysk dla małych obrazów.

\subsection{Zestawienie wersji}

Tabela~\ref{tab:wersje} zestawia, co zmienia każda wersja względem W1, i jakie czynniki będą ograniczać jej przyspieszenie.

\begin{table}[ht]
	\centering
	\caption{Wersje programu: rodzaj przyspieszenia i oczekiwane ograniczenia.}
	\label{tab:wersje}
	\small
	\begin{tabularx}{\textwidth}{@{}lLLL@{}}
		\toprule
		Wersja & Co zmienia względem W1 & Poziom równoległości \cite{czech} & Oczekiwane ograniczenia \\
		\midrule
		W1 & -- & brak & -- \\
		W2 & brzegi poza pętlą, maska w czasie kompilacji, splot rozdzielny, \kod{int16}, bufory wielokrotnego użytku & brak (1 rdzeń) & niewielki zysk rozdzielności dla $k = 3$; drugi przebieg skalowania \\
		W3 & 8 pikseli na instrukcję (AVX2), skalowanie wektorowe & poziom instrukcji (SIMD) & rozszerzanie bajtów, brzegi skalarne, pamięć \\
		W4 & podział wierszy pomiędzy $N$ wątków, pula, scheduler & poziom zadań (wątki, rdzenie) & część sekwencyjna, bariera, SMT, przepustowość pamięci \\
		W5 & obliczenia na karcie graficznej, 1 wątek na piksel & masowa równoległość danych & przesył PCIe, podwójna precyzja \\
		\bottomrule
	\end{tabularx}
\end{table}

\subsection{Propozycje dalszych przyspieszeń z pomocą AI}

Po wykonaniu pomiarów model językowy (Claude), mając dostęp do kodu i wyników, zaproponuje dalsze ulepszenia, których projekt jeszcze nie zawiera (np. wektorowe skalowanie we wszystkich wersjach, instrukcje AVX-512, duże strony pamięci, aktywne oczekiwanie wątków puli, równoległe kodowanie pliku wynikowego, nakładanie przesyłu i obliczeń na GPU). Propozycje zostaną opisane wraz z oczekiwanym zyskiem i uzasadnieniem, a wybrane -- sprawdzone pomiarem. W sprawozdaniu zostanie zaznaczone, że propozycje pochodzą od tego samego narzędzia, które wspierało implementację.

%======================================================================
\section{Metoda, technologie i narzędzia implementacji}

Program zostanie napisany w języku C++ (standard C++17) z wykorzystaniem biblioteki STL, w szczególności mechanizmów wielowątkowości: \kod{std::thread}, \kod{std::mutex}, \kod{std::condition\_variable} i \kod{std::atomic} \cite{williams,cppreference}. Wersja W3 użyje funkcji wewnętrznych kompilatora dla rozszerzenia AVX2 (\kod{immintrin.h}) \cite{intel}, kompilowanych tylko dla funkcji tej wersji, aby pozostała część programu działała na każdym procesorze x86-64. Wersja W5 użyje OpenCL~1.2 ze sterownikiem karty NVIDIA. Odczyt i zapis plików graficznych zrealizują jednoplikowe biblioteki \kod{stb\_image} i \kod{stb\_image\_write} \cite{stb}.

Logika obliczeń zostanie zamknięta w bibliotece \kod{edge\_core}, w której wszystkie wersje są wywoływane przez jedną funkcję. Gwarantuje to identyczne granice pomiaru czasu dla wszystkich wersji. Pomiary wykona konsolowa aplikacja \kod{edge\_bench} oparta na module Qt~Core \cite{qt}, uruchamiana z pliku opisującego eksperyment (JSON) i zapisująca surowe wyniki do plików CSV wraz z opisem warunków pomiaru (procesor, kompilator, flagi, wersja kodu). Projekt będzie budowany narzędziem CMake kompilatorem GCC~13 w trybie Release (\kod{-O3}). Skrypty w języku Python (biblioteki NumPy, pandas, Matplotlib, Pillow) uruchomią eksperymenty kolejno, przetworzą wyniki i wygenerują wykresy oraz tabele; dokumentacja zostanie przygotowana w systemie \LaTeX.

Pomiary zostaną wykonane na komputerze z procesorem AMD Ryzen 9 9900X (12 rdzeni, 24 wątki sprzętowe, 2 bloki rdzeni po 32~MB pamięci L3, obsługa AVX2) i kartą graficzną NVIDIA GeForce GTX 1050 Ti (4~GB) pod systemem Ubuntu 24.04.

%======================================================================
\section{Sposób testowania i oceny jakości rozwiązań}
\label{sec:testy}

\paragraph{Poprawność.} Wynik każdej wersji zostanie porównany z wynikiem W1 dla tego samego obrazu, maski i sposobu skalowania. Zestaw testowy obejmie przypadki brzegowe: obraz $1 \times 1$, jedną kolumnę, jeden wiersz, obraz mniejszy od maski $7\times7$, mniej wierszy niż wątków i wymiary nieparzyste, a dla W4 i W5 -- różne liczby wątków, strategie przydziału, rozmiary fragmentu i grupy roboczej. Wymagany jest wynik identyczny co do bitu; dla wariantu GPU w pojedynczej precyzji zostanie podana największa różnica.

\paragraph{Pomiary czasu.} Dla każdej wersji zostaną wykonane dwa pomiary:
\begin{enumerate}
	\item \textbf{sam algorytm} -- od obrazu w pamięci do gotowego obrazu wynikowego; dla W5 łącznie z przesłaniem danych na kartę i z powrotem;
	\item \textbf{cały program} -- wczytanie i zdekodowanie pliku obrazu, algorytm, zapis wyniku do pliku.
\end{enumerate}
Analiza skupi się na samym algorytmie, ale dla każdej wersji zostanie pokazany rozkład czasu całego programu na wczytanie obrazu, algorytm i zapis wyniku, ponieważ operacje plikowe są sekwencyjne i -- zgodnie z prawem Amdahla~\eqref{eq:amdahl} -- mogą ograniczać zysk całego programu bardziej niż sam algorytm. Czas algorytmu zostanie dodatkowo rozbity na etapy (splot, skalowanie, a dla W5 przesył do karty, obliczenia i przesył z karty).

Każda konfiguracja zostanie wykonana najpierw 1--2 razy bez zapisu wyniku (rozgrzewka), a następnie 15--21 razy. Jako wartość zostanie przyjęta mediana, a jako rozrzut -- zakres od minimum do maksimum \cite{hoefler}. Czas będzie mierzony zegarem \kod{std::chrono::steady\_clock}. Kolejność konfiguracji zostanie wylosowana ze stałym ziarnem, a pomiar uruchomiony tylko wtedy, gdy inne procesy zajmują mniej niż 5\,\% czasu procesora; regulator częstotliwości procesora zostanie ustawiony w tryb \emph{performance}. Przyspieszenie zawsze zostanie podane razem z czasem bezwzględnym wersji bazowej.

\paragraph{Dane testowe.} Zdjęcie NASA Blue Marble ($21\,600 \times 10\,800$ pikseli) przeskalowane do 14 rozmiarów od $256 \times 144$ do $10\,240 \times 5\,760$ (0,04--59~Mpix), obrazy o tej samej liczbie pikseli i różnych proporcjach, zdjęcia z zestawu Kodak, obraz Lena oraz obrazy syntetyczne.

\paragraph{Plan eksperymentów.} Porównanie wersji W1--W5 dla obrazów HD, 4K i 8K i masek $3\times3$ i $7\times7$ (sam algorytm i cały program); wpływ rozmiaru maski na każdą wersję; dla W4 -- liczba wątków ($S(N)$, $E(N)$, dopasowanie prawa Amdahla), strategia przydziału i rozmiar fragmentu, sposób wyznaczania maksimum, pula w porównaniu z nowymi wątkami, przypisanie do rdzeni, cały program w funkcji liczby wątków; dla W5 -- rozkład czasu na przesył i obliczenia, liczba wątków i rozmiar grupy roboczej, precyzja; dla wszystkich wersji -- czas i przepustowość w funkcji rozmiaru obrazu oraz rozmiar obrazu, od którego dana wersja jest szybsza od W1.

\paragraph{Jakość wykrywania krawędzi.} Zostanie sprawdzone, czy większa maska daje lepszy wynik. Na syntetycznym obrazie z figurami o stałej jasności, dla którego położenie krawędzi jest znane co do piksela, zostanie dodany szum gaussowski o odchyleniu standardowym $\sigma$ od 0 do 40 poziomów szarości. Wynik każdej maski zostanie zamieniony na obraz binarny progowaniem (próg dobierany osobno dla każdej maski i poziomu szumu tak, aby dawał najlepszy wynik, co porównuje maski w ich najlepszym ustawieniu) i porównany z prawdziwymi krawędziami miarą Pratta \cite{pratt}:
\begin{equation}
	F = \frac{1}{\max(N_I, N_D)} \sum_{j=1}^{N_D} \frac{1}{1 + \alpha\, d_j^{2}}, \qquad \alpha = \tfrac{1}{9},
\end{equation}
gdzie $N_I$ -- liczba prawdziwych pikseli krawędzi, $N_D$ -- liczba wykrytych, $d_j$ -- odległość $j$-tego wykrytego piksela od najbliższej prawdziwej krawędzi. $F = 1$ oznacza wykrycie idealne; miara karze zarówno brakujące i fałszywe krawędzie, jak i krawędzie przesunięte lub pogrubione. Oczekuje się, że większa maska, która mocniej wygładza obraz, będzie odporniejsza na szum, ale da grubsze i gorzej położone krawędzie przy obrazie bez szumu.

\paragraph{Prezentacja wyników.} Wyniki zostaną przedstawione na wykresach generowanych bezpośrednio z surowych danych: czas i przyspieszenie każdej wersji względem W1 (z czasem bezwzględnym W1), rozkład czasu na etapy dla każdej wersji, przyspieszenie $S(N)$ wersji W4 z linią idealną $S = N$ i krzywą prawa Amdahla, czas w funkcji rozmiaru obrazu, jakość wykrywania w funkcji szumu. Wnioski zostaną poparte liczbami i odniesione do mechanizmów: liczby operacji, szerokości wektora, części sekwencyjnej, przepustowości pamięci i topologii procesora.

%======================================================================
\begin{thebibliography}{99}
	\bibitem{gonzalez}
	Gonzalez R.~C., Woods R.~E., \emph{Digital Image Processing}, 4th ed., Pearson, 2018.

	\bibitem{tadeusiewicz}
	Tadeusiewicz R., Korohoda P., \emph{Komputerowa analiza i przetwarzanie obrazów}, Wydawnictwo Fundacji Postępu Telekomunikacji, Kraków, 1997.

	\bibitem{czech}
	Czech Z., \emph{Wprowadzenie do obliczeń równoległych}, PWN, Warszawa, 2010.

	\bibitem{williams}
	Williams A., \emph{C++ Concurrency in Action}, 2nd ed., Manning, 2019.

	\bibitem{openmp}
	OpenMP Architecture Review Board, \emph{OpenMP Application Programming Interface}, Version 5.2, 2021.

	\bibitem{amdahl}
	Amdahl G.~M., \emph{Validity of the single processor approach to achieving large scale computing capabilities}, AFIPS Conference Proceedings, vol.~30, 1967.

	\bibitem{drepper}
	Drepper U., \emph{What Every Programmer Should Know About Memory}, Red Hat, 2007.

	\bibitem{intel}
	Intel Intrinsics Guide, \url{https://www.intel.com/content/www/us/en/docs/intrinsics-guide/}, 09.11.2026.

	\bibitem{opencl}
	Khronos Group, \emph{The OpenCL Specification}, Version 1.2, 2012.

	\bibitem{pratt}
	Abdou I.~E., Pratt W.~K., \emph{Quantitative design and evaluation of enhancement/thresholding edge detectors}, Proceedings of the IEEE, vol.~67, no.~5, 1979, s.~753--763.

	\bibitem{hoefler}
	Hoefler T., Belli R., \emph{Scientific Benchmarking of Parallel Computing Systems: Twelve Ways to Tell the Masses when Reporting Performance Results}, Proceedings of SC'15, ACM, 2015.

	\bibitem{cppreference}
	\url{https://en.cppreference.com/w/cpp/thread}, 19.10.2026.

	\bibitem{qt}
	Qt 6 Documentation, \url{https://doc.qt.io/qt-6/}, 09.11.2026.

	\bibitem{stb}
	\url{https://github.com/nothings/stb}, 19.10.2026.
\end{thebibliography}

\end{document}
